% Zumdahl Chapter 5 test -- 20 MCQ + 2 FRQ. 50 min.
\documentclass{apchem-exam}

\apcunitword{Chapter}
\apcunit{5}
\apcunittitle{Gases}
\apcdoctitle{Chapter Test}
\apcsubtitle{Zumdahl \S5.1--5.9 \textbullet\ 50 minutes \textbullet\ periodic table and constants provided}

\begin{document}
\apctitleblock

\examsection{I}{Multiple Choice}{20 questions \textbullet\ 30 minutes \textbullet\ 1 point each}{\nocalculator}

% ---- pressure and empirical laws -------------------------------------------
\begin{mcq}
  Which pressure is equivalent to \SI{2.00}{\atm}?
  \choices
    \choice{\SI{380}{\torr}}
    \correct{\SI{1520}{\torr}}
    \choice{\SI{760}{\torr}}
    \choice{\SI{101}{\kilo\pascal}}
\end{mcq}
\why{$2.00 \times 760 = 1520$.}
\distractor{D}{that is 1 atm in kPa}

\begin{mcq}
  A sealed can collapses when the steam inside condenses because
  \choices
    \choice{a vacuum pulls the walls inward}
    \correct{the internal pressure drops far below atmospheric pressure}
    \choice{condensation makes the metal contract}
    \choice{the water expands as it cools}
\end{mcq}
\why{Atmospheric pressure was always present; the inside simply stopped
pushing back.}
\distractor{A}{vacuums do not pull --- pressure differences push}

\begin{mcq}
  A gas at \SI{4.0}{\atm} in \SI{3.0}{\liter} is expanded to
  \SI{12.0}{\liter} at constant temperature. Its new pressure is
  \choices
    \choice{\SI{16}{\atm}}
    \correct{\SI{1.0}{\atm}}
    \choice{\SI{4.0}{\atm}}
    \choice{\SI{0.25}{\atm}}
\end{mcq}
\why{Boyle: $(4.0)(3.0)/12.0 = 1.0$.}

\begin{mcq}
  A gas occupies \SI{2.0}{\liter} at \SI{200}{\kelvin}. At constant
  pressure, its volume at \SI{400}{\kelvin} is
  \choices
    \choice{\SI{1.0}{\liter}}
    \correct{\SI{4.0}{\liter}}
    \choice{\SI{2.0}{\liter}}
    \choice{\SI{8.0}{\liter}}
\end{mcq}
\why{Charles: absolute temperature doubles, so volume doubles.}

\begin{mcq}
  Which variable must always be expressed in kelvin when using a gas law?
  \choices
    \choice{pressure}
    \choice{volume}
    \correct{temperature}
    \choice{moles}
\end{mcq}
\why{Celsius has an arbitrary zero, so ratios of Celsius temperatures are
meaningless.}

% ---- ideal gas law ---------------------------------------------------------
\begin{mcq}
  What volume does \SI{2.00}{\mole} of an ideal gas occupy at STP?
  \choices
    \choice{\SI{22.4}{\liter}}
    \correct{\SI{44.8}{\liter}}
    \choice{\SI{11.2}{\liter}}
    \choice{\SI{2.00}{\liter}}
\end{mcq}
\why{$2.00 \times 22.42$.}

\begin{mcq}
  The molar volume of \SI{22.4}{\liter\per\mole} may be used only when
  \choices
    \choice{the gas is diatomic}
    \choice{the gas is at room temperature}
    \correct{the gas is at \SI{0}{\celsius} and \SI{1}{\atm}}
    \choice{the gas behaves non-ideally}
\end{mcq}
\why{It is derived from $PV = nRT$ at STP specifically.}

\begin{mcq}
  Which expression gives the density of an ideal gas?
  \choices
    \correct{$PM/RT$}
    \choice{$RT/PM$}
    \choice{$MRT/P$}
    \choice{$nRT/V$}
\end{mcq}
\why{From $PV = (m/M)RT$, rearranged for $m/V$.}

\begin{mcq}
  At constant temperature and pressure, which gas has the greatest density?
  \choices
    \choice{\ce{H2}}
    \choice{\ce{He}}
    \choice{\ce{N2}}
    \correct{\ce{CO2}}
\end{mcq}
\why{At fixed $P$ and $T$, density is proportional to molar mass.}

\begin{mcq}
  A \SI{1.00}{\liter} sample of gas at STP has a mass of \SI{1.96}{\gram}.
  Its molar mass is closest to
  \choices
    \choice{\SI{22.4}{\gram\per\mole}}
    \correct{\SI{44}{\gram\per\mole}}
    \choice{\SI{2.0}{\gram\per\mole}}
    \choice{\SI{88}{\gram\per\mole}}
\end{mcq}
\why{$M = d \times 22.42 = 1.96 \times 22.42 \approx 44$.}

% ---- gas stoichiometry -----------------------------------------------------
\begin{mcq}
  For \ce{N2(g) + 3 H2(g) -> 2 NH3(g)} with all gases at the same $T$ and
  $P$, \SI{9.0}{\liter} of \ce{H2} reacts completely with
  \choices
    \correct{\SI{3.0}{\liter} of \ce{N2}}
    \choice{\SI{9.0}{\liter} of \ce{N2}}
    \choice{\SI{27}{\liter} of \ce{N2}}
    \choice{\SI{6.0}{\liter} of \ce{N2}}
\end{mcq}
\why{Volumes combine in the coefficient ratio 1:3 at fixed $T$ and $P$.}

\begin{mcq}
  \ce{CaCO3(s) -> CaO(s) + CO2(g)}. Decomposing \SI{1.00}{\mole} of
  \ce{CaCO3} produces what volume of \ce{CO2} at STP?
  \choices
    \choice{\SI{11.2}{\liter}}
    \correct{\SI{22.4}{\liter}}
    \choice{\SI{44.8}{\liter}}
    \choice{\SI{100}{\liter}}
\end{mcq}
\why{1:1 ratio, then one molar volume.}

\begin{mcq}
  Gas volumes may be compared directly to identify a limiting reactant only
  when
  \choices
    \choice{the gases have equal molar masses}
    \correct{all gases are at the same temperature and pressure}
    \choice{the reaction occurs at STP}
    \choice{the gases are all diatomic}
\end{mcq}
\why{Then volume is proportional to moles (Avogadro's principle).}

% ---- partial pressures -----------------------------------------------------
\begin{mcq}
  A flask holds \SI{0.20}{\mole} \ce{He} and \SI{0.60}{\mole} \ce{Ar} at a
  total pressure of \SI{4.0}{\atm}. The partial pressure of helium is
  \choices
    \correct{\SI{1.0}{\atm}}
    \choice{\SI{2.0}{\atm}}
    \choice{\SI{3.0}{\atm}}
    \choice{\SI{0.80}{\atm}}
\end{mcq}
\why{$\chi_{\ce{He}} = 0.20/0.80 = 0.25$; $0.25 \times 4.0 = 1.0$.}

\begin{mcq}
  In a mixture of ideal gases, the partial pressure exerted by one component
  depends on
  \choices
    \choice{its molar mass}
    \choice{the molar masses of the other gases}
    \correct{its mole fraction and the total pressure}
    \choice{the size of its molecules}
\end{mcq}
\why{$P_A = \chi_A P_{\text{total}}$ --- identity is irrelevant.}

\begin{mcq}
  A gas collected over water at \SI{25}{\celsius} shows a total pressure of
  \SI{760}{\torr}. If the vapour pressure of water is \SI{23.8}{\torr}, the
  partial pressure of the collected gas is
  \choices
    \choice{\SI{783.8}{\torr}}
    \correct{\SI{736.2}{\torr}}
    \choice{\SI{760}{\torr}}
    \choice{\SI{23.8}{\torr}}
\end{mcq}
\why{Subtract the water vapour's contribution.}
\distractor{A}{added instead of subtracting --- a very common slip}

\begin{mcq}
  A student collecting gas over water forgets to subtract the water vapour
  pressure. The calculated number of moles of gas will be
  \choices
    \correct{too high}
    \choice{too low}
    \choice{unaffected}
    \choice{negative}
\end{mcq}
\why{Using the inflated total pressure in $n = PV/RT$ overstates $n$.}

% ---- KMT and real gases ----------------------------------------------------
\begin{mcq}
  Two gases at the same temperature must have the same
  \choices
    \choice{average molecular speed}
    \correct{average kinetic energy}
    \choice{density}
    \choice{molar mass}
\end{mcq}
\why{$KE_{\text{avg}}$ depends only on absolute temperature.}

\begin{mcq}
  As the temperature of a gas rises, its Maxwell--Boltzmann distribution
  \choices
    \choice{narrows and shifts left}
    \correct{broadens and shifts right, with unchanged area}
    \choice{keeps its shape but grows in area}
    \choice{becomes a straight line}
\end{mcq}
\why{Area equals the total number of molecules, which does not change.}

\begin{mcq}
  A real gas shows $PV/nRT$ less than 1. The best explanation is that
  \choices
    \correct{intermolecular attractions reduce the force of wall collisions}
    \choice{the molecules occupy significant volume}
    \choice{the temperature is measured in Celsius}
    \choice{the gas is decomposing}
\end{mcq}
\why{Attractions lower the measured pressure below the ideal prediction.}
\distractor{B}{that effect pushes the ratio ABOVE 1}

\examsection{II}{Free Response}{2 questions \textbullet\ 20 minutes}{\calculatorok}

% ---- FRQ 1 -----------------------------------------------------------------
\frq{short}{4}{Zumdahl \S5.4--5.5 \textbullet\ gas stoichiometry and collection}

Magnesium reacts with hydrochloric acid:
\[ \ce{Mg(s) + 2 HCl(aq) -> MgCl2(aq) + H2(g)} \]
A \SI{0.243}{\gram} sample of Mg ($M = \SI{24.31}{\gram\per\mole}$) reacts
completely. The hydrogen is collected over water at \SI{25}{\celsius}, where
the vapour pressure of water is \SI{23.8}{\torr}. The total pressure is
\SI{752}{\torr}.

\begin{parts}
  \item Calculate the moles of \ce{H2} produced.
        \answerlines[2]{$n(\ce{Mg}) = 0.243/24.31 = \SI{0.0100}{\mole}$;
        1:1 ratio $\to \SI{0.0100}{\mole}$ \ce{H2}.}
  \item Determine the partial pressure of \ce{H2} in atm.
        \answerlines[2]{$752 - 23.8 = \SI{728.2}{\torr}$;
        $\div 760 = \SI{0.9582}{\atm}$}
  \item Calculate the volume of gas collected. \workspace[2.2cm]
        \keyonly{\begin{apcanswer}$V = nRT/P = (0.0100)(0.08206)(298)/0.9582
        = \SI{0.255}{\liter}$\end{apcanswer}}
\end{parts}

\begin{rubric}
  \rpoint{1}{(a) \SI{0.0100}{\mole} with the 1:1 ratio used}
  \rpoint{1}{(b) subtracts the water vapour pressure}
  \rpoint{1}{(b) converts to atm correctly}
  \rpoint{1}{(c) \SI{0.255}{\liter} using kelvin and the corrected pressure}
\end{rubric}

% ---- FRQ 2 -----------------------------------------------------------------
\frq{short}{4}{Zumdahl \S5.6, 5.8 \textbullet\ KMT and real behaviour}

Consider samples of \ce{He} and \ce{CO2}, each at \SI{300}{\kelvin}.

\begin{parts}
  \item Compare their average kinetic energies and their average molecular
        speeds, with justification.
        \answerlines[3]{Average kinetic energies are equal --- $KE$ depends
        only on absolute temperature. Helium moves much faster, because
        $u_{\text{rms}} = \sqrt{3RT/M}$ and helium's molar mass is about
        eleven times smaller.}
  \item At \SI{200}{\kelvin} and \SI{50}{\atm}, which gas is expected to
        deviate more from ideal behaviour? Justify.
        \answerlines[2]{\ce{CO2} --- it is larger and has significantly
        stronger dispersion forces, so both the volume and attraction
        corrections matter more than for tiny, weakly interacting helium.}
  \item State one change in conditions that would make both gases behave
        more ideally, and explain why.
        \answerlines[2]{Raise the temperature (or lower the pressure).
        Faster molecules overcome intermolecular attractions; lower pressure
        keeps particles far apart so their own volume stays negligible.}
\end{parts}

\begin{rubric}
  \rpoint{1}{(a) equal kinetic energies, tied to temperature}
  \rpoint{1}{(a) He faster, with the $1/\sqrt{M}$ reasoning}
  \rpoint{1}{(b) \ce{CO2} with a size or IMF justification}
  \rpoint{1}{(c) names a valid change AND the matching reason}
\end{rubric}

\examtotal

\end{document}
